% ============================================================
%  Constraint and Transformation
%  An Integrated Introduction to Algebra, Geometry,
%  Trigonometry, and Stoichiometry
%
%  Author: Flyxion
%  Second Edition (Revised and Expanded)
% ============================================================

\documentclass[11pt, openany]{book}

% ---- Core packages ------------------------------------------
\usepackage[T1]{fontenc}
\usepackage[utf8]{inputenc}
\usepackage[margin=1.15in, top=1.25in, bottom=1.25in]{geometry}
\usepackage{amsmath, amssymb, amsthm}
\usepackage{mathtools}
\usepackage{xcolor}
\usepackage{tikz}
\usetikzlibrary{arrows.meta, calc, shapes.geometric, positioning, decorations.markings}
\usepackage{pgfplots}
\pgfplotsset{compat=1.18}
\usepackage{booktabs}
\usepackage{array}
\usepackage{enumitem}
\usepackage{tcolorbox}
\tcbuselibrary{theorems, skins, breakable}
\usepackage{fancyhdr}
\usepackage{sectsty}
\usepackage{hyperref}
\usepackage{graphicx}
\usepackage{multicol}
\usepackage{caption}
\usepackage[numbers, sort&compress]{natbib}

% ---- Color palette ------------------------------------------
\definecolor{ctBlue}{RGB}{25, 90, 150}
\definecolor{ctTeal}{RGB}{0, 130, 130}
\definecolor{ctGold}{RGB}{180, 130, 20}
\definecolor{ctRed}{RGB}{170, 40, 40}
\definecolor{ctGray}{RGB}{90, 90, 95}
\definecolor{ctLightBlue}{RGB}{220, 235, 248}
\definecolor{ctLightTeal}{RGB}{215, 240, 238}
\definecolor{ctLightGold}{RGB}{252, 245, 215}
\definecolor{ctLightRed}{RGB}{250, 225, 222}
\definecolor{ctLightGray}{RGB}{245, 245, 247}
\definecolor{ctPurple}{RGB}{90, 40, 120}
\definecolor{ctLightPurple}{RGB}{235, 220, 248}

% ---- Section fonts ------------------------------------------
\chapterfont{\color{ctBlue}}
\sectionfont{\color{ctBlue}}
\subsectionfont{\color{ctGray}}

% ---- Hyperref -----------------------------------------------
\hypersetup{
  colorlinks=true, linkcolor=ctBlue, urlcolor=ctTeal,
  citecolor=ctGray,
  pdftitle={Constraint and Transformation},
  pdfauthor={Flyxion},
}

% ---- Headers / footers --------------------------------------
\pagestyle{fancy}
\fancyhf{}
\fancyhead[LE]{\small\color{ctGray}\textit{Constraint and Transformation}}
\fancyhead[RO]{\small\color{ctGray}\textit{\nouppercase{\leftmark}}}
\fancyfoot[C]{\small\color{ctGray}\thepage}
\renewcommand{\headrulewidth}{0.4pt}

% ============================================================
% THEOREM / PROOF ENVIRONMENTS
% ============================================================
\theoremstyle{plain}
\newtheorem{theorem}{Theorem}[chapter]
\newtheorem{proposition}[theorem]{Proposition}
\newtheorem{lemma}[theorem]{Lemma}
\newtheorem{corollary}[theorem]{Corollary}

\theoremstyle{definition}
\newtheorem{defn}[theorem]{Definition}

\theoremstyle{remark}
\newtheorem{remark}[theorem]{Remark}
\newtheorem{note}[theorem]{Note}

\newenvironment{proofsketch}{\noindent\textit{Proof Sketch.}\;}{\hfill$\square$\par\medskip}

% ============================================================
% TCOLORBOX ENVIRONMENTS
% ============================================================
\tcbset{
  principleStyle/.style={
    enhanced, breakable,
    colback=ctLightBlue, colframe=ctBlue,
    fonttitle=\bfseries\color{white}, colbacktitle=ctBlue,
    attach boxed title to top left={yshift=-2mm, xshift=4mm},
    boxed title style={colback=ctBlue, rounded corners},
    left=4pt, right=4pt, top=8pt, bottom=4pt
  },
  exampleStyle/.style={
    enhanced, breakable,
    colback=ctLightTeal, colframe=ctTeal,
    fonttitle=\bfseries\color{white}, colbacktitle=ctTeal,
    attach boxed title to top left={yshift=-2mm, xshift=4mm},
    boxed title style={colback=ctTeal, rounded corners},
    left=4pt, right=4pt, top=8pt, bottom=4pt
  },
  definitionStyle/.style={
    enhanced, breakable,
    colback=ctLightGold, colframe=ctGold,
    fonttitle=\bfseries\color{white}, colbacktitle=ctGold,
    attach boxed title to top left={yshift=-2mm, xshift=4mm},
    boxed title style={colback=ctGold, rounded corners},
    left=4pt, right=4pt, top=8pt, bottom=4pt
  },
  historyStyle/.style={
    enhanced, breakable,
    colback=ctLightGray, colframe=ctGray,
    fonttitle=\bfseries\color{white}, colbacktitle=ctGray,
    attach boxed title to top left={yshift=-2mm, xshift=4mm},
    boxed title style={colback=ctGray, rounded corners},
    left=4pt, right=4pt, top=8pt, bottom=4pt
  },
  appendixStyle/.style={
    enhanced, breakable,
    colback=ctLightPurple, colframe=ctPurple,
    fonttitle=\bfseries\color{white}, colbacktitle=ctPurple,
    attach boxed title to top left={yshift=-2mm, xshift=4mm},
    boxed title style={colback=ctPurple, rounded corners},
    left=4pt, right=4pt, top=8pt, bottom=4pt
  },
}

\newtcolorbox{principle}[1][]{principleStyle, title={Guiding Principle #1}}
\newtcolorbox{myexample}[1][]{exampleStyle, title={Worked Example #1}}
\newtcolorbox{mydefinition}[1][]{definitionStyle, title={Definition #1}}
\newtcolorbox{history}[1][]{historyStyle, title={Historical Note #1}}
\newtcolorbox{deepnote}[1][]{appendixStyle, title={Deeper Structure #1}}

% ============================================================
% EXERCISE ENVIRONMENTS
% ============================================================
\newcounter{exercise}[chapter]
\renewcommand{\theexercise}{\thechapter.\arabic{exercise}}

\newenvironment{practice}{%
  \refstepcounter{exercise}%
  \par\medskip\noindent
  \textbf{\color{ctTeal}Practice \theexercise.}\enskip}{\par}

\newenvironment{exercise}{%
  \refstepcounter{exercise}%
  \par\medskip\noindent
  \textbf{\color{ctBlue}Exercise \theexercise.}\enskip}{\par}

\newenvironment{challenge}{%
  \refstepcounter{exercise}%
  \par\medskip\noindent
  \textbf{\color{ctGold}Challenge \theexercise.}\enskip}{\par}

\newenvironment{exercises}[1][Exercises]{%
  \section*{#1}\addcontentsline{toc}{section}{#1}}{}

% ============================================================
% TIKZ STYLES
% ============================================================
\tikzset{
  qty/.style={circle, draw=ctGray, fill=ctLightBlue,
              minimum size=7mm, font=\small},
  eqedge/.style={thick, ctBlue},
  arrow/.style={-{Stealth[length=2mm]}, ctBlue, thick},
  varrow/.style={-{Stealth[length=2mm]}, ctTeal, thick},
}

% ============================================================
% MACROS
% ============================================================
\newcommand{\R}{\mathbb{R}}
\newcommand{\N}{\mathbb{N}}
\newcommand{\Z}{\mathbb{Z}}
\newcommand{\Q}{\mathbb{Q}}
\newcommand{\abs}[1]{\left\lvert#1\right\rvert}
\newcommand{\norm}[1]{\left\lVert#1\right\rVert}
\newcommand{\pd}[2]{\frac{\partial #1}{\partial #2}}
\newcommand{\mol}{{\,\text{mol}}}
\newcommand{\gpm}{{\,\text{g/mol}}}
\DeclareMathOperator{\Hom}{Hom}
\DeclareMathOperator{\Ob}{Ob}
\DeclareMathOperator{\id}{id}

% ============================================================
\begin{document}
\frontmatter

% ---- Title page --------------------------------------------
\begin{titlepage}
\centering\vspace*{1.8cm}
{\color{ctBlue}\rule{\linewidth}{1.5pt}}\\[0.5cm]
{\Huge\bfseries\color{ctBlue}Constraint and Transformation}\\[0.4cm]
{\LARGE\color{ctGray}An Integrated Introduction to}\\[0.2cm]
{\LARGE\bfseries\color{ctTeal}
  Algebra, Geometry, Trigonometry,\\[0.1cm]
  and Chemical Stoichiometry}\\[0.5cm]
{\color{ctBlue}\rule{\linewidth}{1.5pt}}\\[1.8cm]

\begin{tikzpicture}[scale=2.0]
  % Unit circle
  \draw[ctBlue!70, thick] (0,0) circle (1);
  \draw[ctGray!40, thin] (-1.3,0)--(1.3,0);
  \draw[ctGray!40, thin] (0,-1.3)--(0,1.3);
  % Angle and point
  \def\ang{40}
  \coordinate (P) at ({cos(\ang)},{sin(\ang)});
  \draw[ctTeal, very thick] (0,0)--(P);
  \draw[ctRed, thick, dashed] ({cos(\ang)},0)--(P)
    node[midway, right, font=\scriptsize, ctRed]{$\sin\theta$};
  \draw[ctGold, thick, dashed] (0,0)--({cos(\ang)},0)
    node[midway, below, font=\scriptsize, ctGold]{$\cos\theta$};
  \draw[ctGray] (0.18,0) arc (0:\ang:0.18)
    node[midway, right=1pt, font=\scriptsize]{$\theta$};
  \filldraw[ctBlue] (P) circle (1.2pt);
  % Constraint label
  \node[font=\tiny, ctGray] at (0,-1.55)
    {$\sin^2\theta+\cos^2\theta=1$};
  \node[font=\tiny, ctGray, right] at (1.05,0) {$1$};
\end{tikzpicture}

\vfill
{\Large\scshape\color{ctGray}Flyxion}\\[0.25cm]
{\color{ctGold}\rule{5cm}{0.8pt}}
\end{titlepage}

% ---- Copyright ---------------------------------------------
\thispagestyle{empty}\vspace*{\fill}
\noindent\textit{Constraint and Transformation}\\
Copyright \textcopyright{} Flyxion. All rights reserved.\\
Typeset in \LaTeX\ with \texttt{pgfplots} and \texttt{tcolorbox}.
\vspace*{\fill}

% ---- ToC ---------------------------------------------------
\tableofcontents

% ============================================================
% NOTATION TABLE
% ============================================================
\chapter*{Notation and Symbols}
\addcontentsline{toc}{chapter}{Notation and Symbols}

The following symbols are used consistently throughout the text.
A change of meaning is always announced explicitly.

\begin{center}
\renewcommand{\arraystretch}{1.55}
\begin{tabular}{>{\bfseries}p{2.2cm} p{10cm}}
\toprule
Symbol & Meaning \\
\midrule
$\R$         & real number line \\
$\N$         & natural numbers $\{0,1,2,\dots\}$ \\
$x,y,z$      & generic real variables \\
$a,b,c$      & constants or geometric side lengths \\
$f,g,h$      & function names \\
$\theta$     & angle (radians unless noted) \\
$r$          & radius of a circle or sphere \\
$n$          & amount of substance (moles, mol) \\
$M$          & molar mass (g/mol) \\
$m$          & mass (grams or kilograms) \\
$L$          & luminosity of a star (watts) \\
$I$          & radiation intensity (W/m$^2$) \\
$d$          & distance \\
$A$          & area \\
$V$          & volume \\
$\Phi$       & scalar field (RSVP plenum density) \\
$\mathbf{v}$ & vector flow field \\
$S$          & entropy density field \\
$\gamma$     & stoichiometric coefficient vector \\
$N$          & element--species incidence matrix \\
$\mathcal{J}$ & variational functional \\
$\preceq$    & causal priority order on events \\
\bottomrule
\end{tabular}
\end{center}

% ============================================================
% PREFACE
% ============================================================
\chapter*{Preface}
\addcontentsline{toc}{chapter}{Preface}

This book is built on one conviction: algebra, geometry, trigonometry, and
chemistry are not separate disciplines accidentally packaged together by
tradition. They are four dialects of the same language --- the language of
\emph{constraint and transformation}.

Every equation encountered here is a claim about balance. Algebra encodes
symbolic balance among quantities. Geometry encodes spatial balance among
dimensions. Trigonometry encodes rotational and cyclic balance. Stoichiometry
encodes material balance in chemical reactions. To understand any one of these
is to recognize a pattern that recurs across all the others: a system remains
coherent because its constitutive relations must simultaneously be satisfied.

This orientation reflects a structural feature of scientific reasoning.
Quantitative systems are defined by relations that must remain internally
consistent as quantities change. A child discovering that a toy machine follows
a hidden rule, an astrophysicist reading the composition of a star from its
spectrum, a political theorist analyzing how collective power emerges from
coordination, a novelist constructing an ecologically stable imaginary world ---
each is solving, in different vocabulary, the same problem: \emph{what
constraints govern this system, and what transformations preserve them?}

\medskip\noindent\textbf{Intellectual influences.}
The book's perspective has been shaped by several intellectual traditions. The
developmental psychology of \citet{Gopnik2009} shows how causal inference in
children is structurally isomorphic with algebraic constraint satisfaction. The
astrophysics of \citet{PayneGaposchkin1925} demonstrates that geometric and
spectroscopic reasoning can reveal chemical composition at stellar distances.
The political philosophy of \citet{Arendt1958} offers a model in which social
systems are organized by relational structure rather than by fixed hierarchies.
The complexity science of \citet{Juarrero1999} formalizes how constraints
propagate through dynamical systems, while the speculative fiction of
\citet{LeGuin1974} explores the societal consequences of such constraint
structures with rare concreteness. The visual mathematical approach of
\citet{Needham1997,Needham2021} informs the book's emphasis on diagrammatic
reasoning.

\medskip\noindent\textbf{How to read this book.}
Each chapter follows a three-step rhythm: \emph{observation} (a phenomenon in
plain language), \emph{formalization} (the mathematical statement), and
\emph{extension} (the same idea in a new domain). The main text is
self-contained for an introductory reader. Each chapter ends with three tiers of
exercises: \textbf{Practice} (computation), \textbf{Exercise} (conceptual
reasoning), and \textbf{Challenge} (open-ended exploration). The appendices
develop the deeper formal architecture --- field theory, category theory, sheaf
semantics, variational principles, and recursive simulation --- for readers who
wish to see how elementary mathematics extends into contemporary research.

\medskip\noindent\textbf{Prerequisites.}
Arithmetic: fractions, decimals, percentages, signed numbers. No prior
algebra, geometry, or chemistry is assumed.

\mainmatter

% ============================================================
% INTRODUCTION
% ============================================================
\chapter{Mathematics as Structured Relation}

\section{The Central Thesis}

Scientific descriptions can be viewed as networks of relationships rather
than lists of isolated facts. This book proposes that the four subjects it
covers --- algebra, geometry, trigonometry, and chemistry --- are four
languages for a single idea: a system is defined not merely by its components
but by the \emph{relations} those components must satisfy.

\begin{mydefinition}[Constraint]
A \textbf{constraint} is any condition restricting the possible values of one
or more quantities. Formally, given variables $x_1,\dots,x_n$, a constraint is
a relation
\[
  F(x_1,\dots,x_n) = 0.
\]
The \textbf{solution set} is
$\mathcal{C}=\{(x_1,\dots,x_n)\in\R^n \mid F(x_1,\dots,x_n)=0\}.$
\end{mydefinition}

\begin{principle}
Every equation in this book is a claim that two descriptions of the same
situation must agree. Solving the equation means finding the values that live
inside the solution set $\mathcal{C}$.
\end{principle}

\section{Four Languages, One Pattern}

The table below introduces the central motif. Each row names a type of
constraint, the domain where it appears, and a representative equation.

\begin{center}
\renewcommand{\arraystretch}{1.6}
\begin{tabular}{>{\bfseries}p{2.8cm} p{3.2cm} p{5.8cm}}
\toprule
\color{ctBlue}Domain & \color{ctBlue}Constraint type & \color{ctBlue}Canonical equation \\
\midrule
Algebra       & symbolic / numerical   & $2x+3=11$ \\
Geometry      & spatial / dimensional  & $a^2+b^2=c^2$ \\
Trigonometry  & cyclic / angular       & $\sin^2\theta+\cos^2\theta=1$ \\
Stoichiometry & material / atomic      & $2\mathrm{H_2+O_2\to2H_2O}$ \\
\bottomrule
\end{tabular}
\end{center}

In every case the equation is a \emph{conservation statement}: it asserts that
some quantity --- a sum, a squared length, a trigonometric magnitude, or an
atom count --- is preserved under transformation.

\section{Constraint Networks}

A single equation imposes one constraint. Real systems impose many constraints
simultaneously. Representing them as a network reveals their mutual
dependencies.

\begin{center}
\begin{tikzpicture}[node distance=1.8cm]
  \node[qty] (x) {$x$};
  \node[qty, right=of x] (y) {$y$};
  \node[qty, below=of x] (z) {$z$};
  \node[qty, below=of y] (w) {$w$};
  \draw[eqedge] (x)--(y) node[midway,above,font=\scriptsize]{$F_1=0$};
  \draw[eqedge] (x)--(z) node[midway,left,font=\scriptsize]{$F_2=0$};
  \draw[eqedge] (y)--(w) node[midway,right,font=\scriptsize]{$F_3=0$};
  \draw[eqedge] (z)--(w) node[midway,below,font=\scriptsize]{$F_4=0$};
  \draw[eqedge] (y)--(z) node[midway,above right,font=\scriptsize]{$F_5=0$};
\end{tikzpicture}
\end{center}

Each node represents a quantity; each edge represents a constraint relating the
two adjacent quantities. A consistent system is one in which all constraints can
be satisfied simultaneously.

\begin{deepnote}
The general form of a linear constraint system is $A\mathbf{x}=\mathbf{b}$,
where $A$ is an $m\times n$ matrix. Appendix~A develops this idea formally,
including its connections to the RSVP field framework and Spherepop
event-history calculus.
\end{deepnote}

\section{The Book's Arc}

The book progresses from local to global. Part~I develops algebraic constraints
on individual quantities. Part~II extends this to spatial relationships. Part~III
describes cyclic and rotational constraints. Part~IV applies all of this to
chemical systems, which are constraint networks over atomic species. Part~V
synthesizes the four domains into a unified picture of structured systems.

\begin{exercises}
\begin{exercise}
Write down three everyday relationships and express each as an equation
$F(x,y)=0$. What does the solution set mean in each case?
\end{exercise}
\begin{challenge}
Draw a constraint network with five quantities and at least six constraints.
Choose a real system (a recipe, a budget, a simple machine) as your model.
Explain why some constraints are necessary and others may be redundant.
\end{challenge}
\end{exercises}

% ============================================================
% PART I --- ALGEBRA
% ============================================================
\part{Algebra: Relations and Constraint Systems}

\chapter{Numbers, Variables, and Representation}

\section{Numbers as Quantities}

A \textbf{number} is a symbolic stand-in for a measurable quantity: a count, a
length, a temperature. Numbers obey arithmetic: addition, subtraction,
multiplication, division.

\begin{mydefinition}[Variable]
A \textbf{variable} is a letter that stands for an unknown or varying quantity.
An \textbf{expression} is a combination of numbers and variables using
arithmetic operations. An expression does not assert equality; it simply
describes a quantity.
\end{mydefinition}

The expression $3x+5$ has a specific numerical value for each specific value
of $x$. It is not an equation and does not claim anything --- it is simply a
description.

\section{The Real Number Line}

Every real number corresponds to a unique point on an infinite line. Distance
from the origin measures magnitude; direction indicates sign.

\begin{center}
\begin{tikzpicture}[scale=1.1]
  \draw[arrow] (-3.5,0)--(3.8,0) node[right,font=\small]{$\R$};
  \foreach \x/\l in {-3/-3,-2/-2,-1/-1,0/0,1/1,2/2,3/3}
    \draw (\x,0.08)--(\x,-0.08) node[below,font=\scriptsize]{\l};
  \filldraw[ctRed] (2.3,0) circle (2.2pt) node[above,font=\scriptsize,ctRed]{$2.3$};
  \filldraw[ctTeal] (-1.6,0) circle (2.2pt) node[above,font=\scriptsize,ctTeal]{$-1.6$};
\end{tikzpicture}
\end{center}

\section{Proportionality: The Simplest Relation}

\begin{theorem}[Direct Proportionality]
If two quantities $x$ and $y$ satisfy $y=kx$ for a constant $k\ne0$, we say
$y$ is \textbf{directly proportional} to $x$ with proportionality constant $k$.
The graph of this relation is a straight line through the origin.
\end{theorem}

\begin{myexample}
A car travels at a constant speed of 60~km/h. Distance $d$ (km) and time $t$
(h) satisfy
\[
  d = 60\,t.
\]
Checking units: $\text{km} = (\text{km/h})\cdot\text{h}$. \checkmark

The slope $k=60$ is the rate of change: each additional hour yields 60 more
kilometres. For $t=2.5$~h, $d=150$~km.
\end{myexample}

\begin{history}[Alison Gopnik and Causal Inference]
In her research on how children understand the world, Alison Gopnik showed
that even toddlers spontaneously construct algebraic-style rules from
observation \citep{Gopnik2009}. A child who notices that a machine lights up
whenever two specific blocks are inserted together is, in effect, discovering a
proportional or Boolean rule governing the system. Elementary algebra makes
this intuition precise.
\end{history}

\begin{exercises}
\begin{practice}
A recipe uses 3 cups of flour per 2 cups of sugar. Write an equation relating
flour $f$ to sugar $s$. How much sugar is needed for 9 cups of flour?
\end{practice}
\begin{exercise}
A toy machine activates according to $A = 3B + 2$, where $B$ is the number of
blue blocks inserted. If the observed activation value is $A=14$, determine $B$
and interpret the result.
\end{exercise}
\begin{challenge}
Collect five proportional relationships from daily life (speed and distance,
ingredients in a recipe, etc.) and estimate each proportionality constant.
Which relationships are only approximately proportional? What causes the
deviation?
\end{challenge}
\end{exercises}

\chapter{Equations and Balance}

\section{Equations as Statements of Equality}

\begin{mydefinition}[Equation and Solution]
An \textbf{equation} is a statement $A=B$ where $A$ and $B$ are expressions.
A \textbf{solution} is any assignment of values to the variables that makes
the statement true.
\end{mydefinition}

An equation is therefore a constraint: it asserts that exactly one value (or
set of values) satisfies the balance condition.

\subsection{The Balance Model}

\begin{center}
\begin{tikzpicture}[scale=0.9]
  \draw[ctBlue, very thick] (-3.1,0.9)--(3.1,0.9);
  \draw[ctBlue, thick] (0,0.9)--(0,-0.2);
  \draw[ctGray] (-0.3,-0.2)--(0.3,-0.2);
  \draw[ctTeal,thick] (-3.1,0.9)--(-3.1,0.4);
  \draw[ctTeal,thick] (3.1,0.9)--(3.1,0.4);
  \draw[ctTeal,thick] (-3.5,0.4)--(-2.7,0.4);
  \draw[ctTeal,thick] (2.7,0.4)--(3.5,0.4);
  \node at (-3.1,0.1){\small$2x+3$};
  \node at (3.1,0.1){\small$11$};
  \node[ctGray,font=\small] at (0,-0.55){balance};
\end{tikzpicture}
\end{center}

Whatever operation is applied to one side of the equation must be applied
identically to the other --- the scale must remain balanced.

\section{Solving Linear Equations}

\begin{theorem}[Isolation Principle]
A linear equation $ax+b=c$ (with $a\ne0$) has the unique solution
\[
  x = \frac{c-b}{a}.
\]
\end{theorem}

\begin{proofsketch}
Subtract $b$ from both sides: $ax=c-b$. Divide both sides by $a\ne0$.
\end{proofsketch}

\begin{myexample}[Balancing a simple equation]
Solve $2x+3=11$.
\begin{align*}
  2x+3 &= 11 && \text{original constraint}\\
  2x   &= 8  && \text{subtract }3\\
  x    &= 4  && \text{divide by }2
\end{align*}
\textit{Verification:} $2(4)+3=11$. \checkmark
\end{myexample}

\begin{myexample}[Gopnik: Hypothesis testing as equation solving]
A child hypothesises that a toy machine follows the rule $A=3G+Y$, where $G$
and $Y$ are the numbers of green and yellow blocks. She observes $A=8$ and
inserts $Y=2$. Solving for $G$:
\[
  3G+2=8 \implies 3G=6 \implies G=2.
\]
\end{myexample}

\section{Dimensional Analysis}

\begin{principle}
Every valid equation involving physical quantities must be \textbf{dimensionally
consistent}: the units on both sides must be identical.
\end{principle}

\begin{myexample}[Converting units]
Convert 72~km/h to m/s:
\[
  72\;\frac{\text{km}}{\text{h}}
  \times\frac{1000\;\text{m}}{1\;\text{km}}
  \times\frac{1\;\text{h}}{3600\;\text{s}}
  = 20\;\frac{\text{m}}{\text{s}}.
\]
\end{myexample}

\begin{exercises}
\begin{practice}
Solve each equation.\\
(a) $5x-7=18$\quad (b) $\dfrac{x}{3}+4=9$\quad
(c) $-2x+10=0$\quad (d) $3(x-2)=15$.
\end{practice}
\begin{exercise}
A reaction produces 0.5~mol of gas per minute. Write the constraint equation
and determine how many moles are produced in 2.5 hours. Include units at every
step.
\end{exercise}
\begin{challenge}
Arendt \citep{Arendt1958} argues that political power $P$ emerges from
coordination, not mere numbers. Model $P=cn$ where $c$ is coordination
strength and $n$ is the number of participants. If $c=0.2$ initially and the
group triples its coordination to $c=0.6$, by what factor does $P$ change?
What does this say about collective action?
\end{challenge}
\end{exercises}

\chapter{Proportions and Scaling}

\section{Ratios and Proportions}

A \textbf{ratio} compares two quantities. A \textbf{proportion} asserts that
two ratios are equal:
\[
  \frac{a}{b}=\frac{c}{d}.
\]

\begin{myexample}[Stellar mass ratios]
In the hydrogen fusion reaction $4\mathrm{H}\to\mathrm{He}$, the mass ratio of
hydrogen consumed to helium produced is approximately $4:1$ (by nucleon count).
If a stellar region contains $10^{12}$ hydrogen nuclei per second, the
production rate of helium nuclei satisfies
\[
  \frac{n_\mathrm{He}}{1}=\frac{n_\mathrm{H}}{4}
  \implies n_\mathrm{He}=\frac{10^{12}}{4}=2.5\times10^{11}\text{ nuclei/s}.
\]
\end{myexample}

\begin{history}[Cecilia Payne-Gaposchkin]
Cecilia Payne-Gaposchkin (1900--1979) discovered in her 1925 doctoral
dissertation that stars are composed primarily of hydrogen and helium --- a
result initially rejected by the astronomical establishment but later
recognised as one of the most important discoveries in astrophysics
\citep{PayneGaposchkin1925}. Her analysis relied on proportional reasoning
connecting spectral line intensities to elemental abundances.
\end{history}

\section{Scaling Laws}

\begin{theorem}[Power Scaling]
If quantity $y$ scales as the $p$-th power of quantity $x$, then $y=kx^p$.
When $x$ is multiplied by a factor $\lambda$, $y$ is multiplied by $\lambda^p$.
\end{theorem}

\begin{myexample}[Le Guin: island resource capacity]
In a circular island settlement of radius $r=2$~km, the total area is
$A=\pi r^2=4\pi\approx12.6\;\text{km}^2$. If agricultural land occupies
half the island and each km$^2$ supports 500 people:
\[
  \text{capacity} = 500\times\frac{4\pi}{2}\times100
  = 500\times200\pi\approx314{,}159\text{ people}.
\]
Doubling the radius to $r=4$~km quadruples the area and thus quadruples the
capacity. Scaling is not linear \citep{LeGuin1974}.
\end{myexample}

\begin{exercises}
\begin{practice}
A car uses 8~L per 100~km. How far can it travel on 22~L?
\end{practice}
\begin{exercise}
A solution is prepared at 5~g of salt per 200~mL. How much salt is needed
for 500~mL? Write the proportion and solve.
\end{exercise}
\begin{challenge}
Le Guin's utopian novel \textit{The Dispossessed} \citep{LeGuin1974} depicts
a society on an arid moon that must balance population against scarce
resources. Suppose the moon's habitable area is $A$~km$^2$, each km$^2$
supports at most $p$ people, and the society aims to keep population at $80\%$
of capacity. Express the target population as a function of $A$ and $p$.
How does the target change if a drought reduces $A$ by $15\%$?
\end{challenge}
\end{exercises}

\chapter{Systems of Equations}

\section{Multiple Constraints, Multiple Unknowns}

When a situation involves two or more unknowns, a single equation constrains
but does not determine. A system of $n$ equations in $n$ unknowns typically
yields a unique solution.

\begin{mydefinition}[System of Equations]
A \textbf{system of $m$ equations in $n$ unknowns} is a collection
$\{F_1=0,\dots,F_m=0\}$ that must all hold simultaneously. The
\textbf{solution} is any tuple satisfying every equation at once.
\end{mydefinition}

\section{Geometric Interpretation}

Each linear equation $ax+by=c$ defines a \textbf{line} in the $xy$-plane. A
system of two such equations corresponds to two lines; the solution is their
intersection.

\begin{center}
\begin{tikzpicture}
\begin{axis}[
  width=7cm, height=6cm, xlabel={$x$}, ylabel={$y$},
  xmin=0,xmax=8, ymin=-1,ymax=10,
  axis lines=center, grid=both, grid style={ctGray!20},
  tick label style={font=\small},
]
  \addplot[ctBlue,thick,domain=0:8]{10-x}
    node[above right,font=\small]{$x+y=10$};
  \addplot[ctTeal,thick,domain=0:8]{3*x-2}
    node[below right,font=\small]{$3x-y=2$};
  \addplot[ctRed, only marks, mark=*, mark size=2.5pt]
    coordinates{(3,7)};
  \node[font=\small,ctRed,above right] at (axis cs:3,7){$(3,7)$};
\end{axis}
\end{tikzpicture}
\end{center}

\begin{myexample}[Juarrero: constraint propagation through a system]
Juarrero \citep{Juarrero1999} emphasises that constraints propagate through
systems: a change in one relation alters the whole. Consider a simple model:
\[
  \begin{cases}
    F = 2L \\
    E = F + L,
  \end{cases}
\]
with $L=100$. Then $F=200$, $E=300$. If a drought reduces $L$ to 80, then
$F=160$ and $E=240$. The constraint cascade propagates through the entire
system.
\end{myexample}

\begin{myexample}[Solving by elimination]
Solve $\{x+y=10,\;3x-y=2\}$.

Adding the two equations: $4x=12$, so $x=3$. Then $3+y=10$, so $y=7$.

\textit{Verification:} $3+7=10$ \checkmark;\; $9-7=2$ \checkmark.
\end{myexample}

\begin{exercises}
\begin{practice}
Solve the system $\{2x+3y=12,\; x-y=1\}$ by both substitution and elimination.
\end{practice}
\begin{exercise}
Interpret geometrically: (a) a system with no solution; (b) a system with
infinitely many solutions. What does each case mean for the constraint
network?
\end{exercise}
\begin{challenge}
A chemist mixes Solution~A (20\% acid) and Solution~B (50\% acid) to produce
300~mL of 35\% acid. Set up and solve the system. Connect this to Juarrero's
idea of constraints propagating through a mixture network.
\end{challenge}
\end{exercises}

\chapter{Matrices and Linear Transformations}

\section{Systems in Matrix Form}

A system $\{a_1x+b_1y=c_1,\;a_2x+b_2y=c_2\}$ may be written as
\[
  A\mathbf{x}=\mathbf{b},
  \qquad
  A=\begin{pmatrix}a_1&b_1\\a_2&b_2\end{pmatrix},\quad
  \mathbf{x}=\begin{pmatrix}x\\y\end{pmatrix},\quad
  \mathbf{b}=\begin{pmatrix}c_1\\c_2\end{pmatrix}.
\]

\begin{mydefinition}[Matrix]
An $m\times n$ \textbf{matrix} is a rectangular array of $mn$ numbers
arranged in $m$ rows and $n$ columns. It encodes a linear transformation
$T:\R^n\to\R^m$ by $T(\mathbf{x})=A\mathbf{x}$.
\end{mydefinition}

\section{Matrices as Transformations of Space}

A $2\times2$ matrix $A$ transforms every point $(x,y)$ in the plane. The
unit square is sent to a parallelogram.

\begin{center}
\begin{tikzpicture}[scale=1.1]
  % Original square
  \fill[ctLightBlue] (0,0)--(1,0)--(1,1)--(0,1)--cycle;
  \draw[ctBlue,thick] (0,0)--(1,0)--(1,1)--(0,1)--cycle;
  \node[font=\scriptsize,ctBlue] at (0.5,0.5){unit};
  % Transformed (shear A = [[1,1],[0,1]])
  \begin{scope}[xshift=3.5cm]
    \fill[ctLightTeal] (0,0)--(1,0)--(2,1)--(1,1)--cycle;
    \draw[ctTeal,thick] (0,0)--(1,0)--(2,1)--(1,1)--cycle;
    \node[font=\scriptsize,ctTeal] at (1,0.5){$A\cdot\square$};
  \end{scope}
  % Arrow
  \draw[arrow] (1.3,0.5)--(3.2,0.5)
    node[midway,above,font=\scriptsize]{$A=\!\scriptsize\begin{pmatrix}1&1\\0&1\end{pmatrix}$};
\end{tikzpicture}
\end{center}

\section{The Stoichiometric Matrix}

Chemical reactions are encoded by a stoichiometric matrix $N$ whose rows
index elements and whose columns index species.

\begin{myexample}[Water formation as matrix equation]
For $2\mathrm{H_2}+\mathrm{O_2}\to2\mathrm{H_2O}$, writing
$\gamma=(-2,-1,2)^\top$ (negative for reactants):
\[
  N=\begin{pmatrix}2&0&2\\0&2&1\end{pmatrix},\qquad
  N\gamma=\begin{pmatrix}2(-2)+0(-1)+2(2)\\0(-2)+2(-1)+1(2)\end{pmatrix}
  =\begin{pmatrix}0\\0\end{pmatrix}.
\]
Conservation is encoded by $N\gamma=\mathbf{0}$.
\end{myexample}

\begin{deepnote}
Appendix~D develops stoichiometry as conservation algebra in full generality,
including the reaction rate equation $\dot{c}=Sr(c)$.
\end{deepnote}

\begin{exercises}
\begin{practice}
Multiply $\begin{pmatrix}2&1\\0&3\end{pmatrix}\begin{pmatrix}4\\-1\end{pmatrix}$.
\end{practice}
\begin{exercise}
Write the system $\{x+2y=7,\;3x-y=5\}$ as $A\mathbf{x}=\mathbf{b}$ and solve
using Cramer's rule or row reduction.
\end{exercise}
\begin{challenge}
Model Arendt's power equation $P=C\mathbf{n}$ as a matrix transformation,
where $\mathbf{n}$ is a vector of participants across three groups and $C$
is a $1\times3$ coordination matrix. What does the rank of $C$ represent?
\end{challenge}
\end{exercises}

% ============================================================
% PART II --- GEOMETRY
% ============================================================
\part{Geometry: Spatial Structure and Invariance}

\chapter{Points, Lines, and Distance}

\section{The Coordinate Plane}

René Descartes' insight was that geometric objects can be described
algebraically. A point is an ordered pair $(x,y)$; a line is the solution
set of a linear equation; a circle is the solution set of a quadratic relation.

\begin{theorem}[Distance Formula]
The distance between points $(x_1,y_1)$ and $(x_2,y_2)$ in the plane is
\[
  d = \sqrt{(x_2-x_1)^2+(y_2-y_1)^2}.
\]
This formula is a consequence of the Pythagorean theorem, which we derive in
the next chapter.
\end{theorem}

\section{Angles and Their Measure}

\begin{mydefinition}[Radian]
One \textbf{radian} is the angle subtended at the centre of a circle by an
arc equal in length to the radius. Since the full circumference is $2\pi r$,
a full rotation equals $2\pi$ radians.
\end{mydefinition}

Conversion: $\theta_\text{rad}=\frac{\pi}{180}\theta_\text{deg}$.

\begin{exercises}
\begin{practice}
Convert: (a) $60°$ to radians; (b) $\frac{3\pi}{4}$ to degrees.
\end{practice}
\begin{exercise}
Show algebraically that the midpoint of the segment from $(x_1,y_1)$ to
$(x_2,y_2)$ is $\bigl(\frac{x_1+x_2}{2},\frac{y_1+y_2}{2}\bigr)$.
\end{exercise}
\end{exercises}

\chapter{Triangles and the Pythagorean Relation}

\section{Right Triangles}

\begin{theorem}[Pythagorean Theorem]
In any right triangle with legs $a,b$ and hypotenuse $c$,
\[
  a^2+b^2=c^2.
\]
The length of the hypotenuse is completely determined by the two legs.
\end{theorem}

\begin{proofsketch}
Place four congruent right triangles inside a square of side $a+b$.
The uncovered central region is a square of side $c$.

\begin{center}
\begin{tikzpicture}[scale=1.5]
  \def\a{1.1}\def\b{0.7}
  \pgfmathsetmacro{\s}{\a+\b}
  \draw[ctGray,thick] (0,0)--(\s,0)--(\s,\s)--(0,\s)--cycle;
  \filldraw[ctLightBlue,draw=ctBlue]
    (0,0)--(\a,0)--(0,\b)--cycle;
  \filldraw[ctLightBlue,draw=ctBlue]
    (\s,0)--(\s,\a)--(\b,0)--cycle;
  \filldraw[ctLightBlue,draw=ctBlue]
    (\s,\s)--(0,\s)--(\s,\s-\b)--cycle;
  \filldraw[ctLightBlue,draw=ctBlue]
    (0,\s)--(0,\s-\a)--(\a,\s)--cycle;
  \draw[ctRed,very thick]
    (\a,0)--(\s,\b)--(\b,\s)--(0,\s-\a)--cycle;
  \node[font=\small,ctRed] at ({\s/2},{\s/2}){$c^2$};
  \node[font=\scriptsize,below] at (\a/2,0){$a$};
  \node[font=\scriptsize,below] at ({\a+\b/2},0){$b$};
\end{tikzpicture}
\end{center}
Area of outer square: $(a+b)^2=a^2+2ab+b^2$. Four triangles: $4\cdot\tfrac12ab=2ab$.
Inner square: $c^2=(a+b)^2-2ab=a^2+b^2$.
\end{proofsketch}

\section{Applications}

\begin{myexample}[Payne-Gaposchkin: radiation and distance]
Light from a star of luminosity $L$ (W) spreads uniformly over a sphere of
radius $r$. Surface area $A=4\pi r^2$ (m$^2$), so intensity is
\[
  I = \frac{L}{4\pi r^2}\;\text{W/m}^2.
\]
For the Sun: $L=3.8\times10^{26}$~W at $r=1.5\times10^{11}$~m:
\[
  4\pi r^2\approx 2.83\times10^{23}\text{ m}^2,\qquad
  I\approx1340\;\text{W/m}^2.
\]
The inverse-square law is a constraint on spatial geometry.
\end{myexample}

\begin{exercises}
\begin{practice}
Find the diagonal of a $7\times24$~cm rectangle.
\end{practice}
\begin{exercise}
Two ships leave the same port. One travels 40~km due north; the other 30~km
due east. How far apart are they? Draw and label the constraint diagram.
\end{exercise}
\begin{challenge}
Star~A has luminosity $2L$ and is at distance $3d$. Star~B has luminosity $L$
and is at distance $d$. Which appears brighter to an observer, and by what
factor? Generalise: express the apparent brightness ratio as a function of
the luminosity and distance ratios.
\end{challenge}
\end{exercises}

\chapter{Similarity and Scaling in Geometry}

\section{Similar Figures}

\begin{mydefinition}[Similarity]
Two figures are \textbf{similar} (written $\sim$) if one can be obtained
from the other by scaling, rotation, or reflection. Corresponding angles are
equal; corresponding lengths are proportional with scale factor $k$.
\end{mydefinition}

\begin{theorem}[Scaling Laws]
If two similar figures have linear scale factor $k$, then areas scale as
$k^2$ and volumes scale as $k^3$.
\end{theorem}

\begin{proofsketch}
Area involves two length dimensions; volume involves three. Each scales as a
power of $k$ equal to the number of dimensions.
\end{proofsketch}

\begin{exercises}
\begin{practice}
Two similar triangles have scale factor $3:5$. If the smaller has area
$27\;\text{cm}^2$, find the area of the larger.
\end{practice}
\begin{exercise}
A spherical star with radius $r$ has surface area $4\pi r^2$. If the star
expands to radius $2r$, by what factor does its surface area increase? Its
volume?
\end{exercise}
\end{exercises}

\chapter{Transformations of the Plane}

\section{Isometries: Transformations Preserving Distance}

A \textbf{transformation} is a function mapping the plane to itself. An
\textbf{isometry} preserves distances:
$d(T(P),T(Q))=d(P,Q)$ for all points $P,Q$.

The three basic isometries are:
\begin{itemize}
\item \textbf{Translation} by vector $\mathbf{v}$: $T(P)=P+\mathbf{v}$.
\item \textbf{Rotation} by angle $\theta$ around origin: encoded by $R(\theta)$.
\item \textbf{Reflection} across a line $\ell$.
\end{itemize}

\begin{center}
\begin{tikzpicture}[scale=0.85]
  % Original triangle
  \coordinate (A) at (0,0); \coordinate (B) at (1.2,0); \coordinate (C) at (0.6,1);
  \filldraw[ctLightBlue,draw=ctBlue] (A)--(B)--(C)--cycle;
  \node[font=\scriptsize,ctBlue] at (0.6,0.3){$T$};
  % Translated
  \begin{scope}[xshift=2.5cm]
    \coordinate (A') at (0,0); \coordinate (B') at (1.2,0); \coordinate (C') at (0.6,1);
    \filldraw[ctLightTeal,draw=ctTeal] (A')--(B')--(C')--cycle;
    \node[font=\scriptsize,ctTeal] at (0.6,0.3){$T'$};
  \end{scope}
  \draw[arrow] (1.3,0.5)--(2.4,0.5) node[midway,above,font=\scriptsize]{translate};
  % Rotated
  \begin{scope}[xshift=5.5cm, rotate=35]
    \coordinate (A'') at (0,0); \coordinate (B'') at (1.2,0); \coordinate (C'') at (0.6,1);
    \filldraw[ctLightGold,draw=ctGold] (A'')--(B'')--(C'')--cycle;
  \end{scope}
  \draw[arrow] (3.9,0.5)--(5.2,0.5) node[midway,above,font=\scriptsize]{rotate};
\end{tikzpicture}
\end{center}

\begin{principle}
Every isometry preserves shape and size. It is a constraint-preserving
transformation: the distance relation $d(P,Q)=c$ becomes $d(T(P),T(Q))=c$.
\end{principle}

\begin{history}[Hannah Arendt and Invariance]
Arendt's analysis of political action distinguishes between acts that alter
the social space and those that maintain it. An isometry is the spatial analogue
of a norm-preserving act: the shape of relations is maintained even as positions
change. The rotation matrix $R(\theta)$ is the canonical mathematical example.
\end{history}

\begin{exercises}
\begin{exercise}
Show that the composition of two rotations around the origin is another
rotation. What is the resulting angle?
\end{exercise}
\begin{challenge}
A rotation by $\theta$ is represented by the matrix
$R(\theta)=\begin{pmatrix}\cos\theta&-\sin\theta\\\sin\theta&\cos\theta\end{pmatrix}$.
Compute $R(\theta)^\top R(\theta)$ and explain why this confirms that $R(\theta)$
is an isometry.
\end{challenge}
\end{exercises}

\chapter{Area, Volume, and Spatial Distribution}

\section{Area Formulas as Algebraic Constraints}

\begin{center}
\renewcommand{\arraystretch}{1.5}
\begin{tabular}{llp{5cm}}
\toprule
Shape & Formula & Constraint expressed \\
\midrule
Rectangle & $A=lw$           & product of perpendicular dimensions \\
Triangle   & $A=\tfrac12 bh$ & half the enclosing rectangle \\
Circle     & $A=\pi r^2$     & area-radius relation \\
Sphere     & $A=4\pi r^2$   & surface of constant curvature \\
Cylinder   & $V=\pi r^2 h$  & circular cross-section times height \\
\bottomrule
\end{tabular}
\end{center}

\begin{myexample}[Le Guin: sustainable population]
A circular settlement of radius $r=3$~km devotes $60\%$ of its area to crops.
Each km$^2$ supports 400 people.
\[
  A = \pi(3)^2 = 9\pi\;\text{km}^2,\qquad
  A_c = 0.6\times9\pi = 5.4\pi\;\text{km}^2.
\]
Population capacity: $400\times5.4\pi\times100\approx679{,}000$.
(Factor of 100 converts km$^2$ to hectares.)
Le Guin's ecological societies are governed by exactly such spatial constraints
\citep{LeGuin1974}.
\end{myexample}

\begin{exercises}
\begin{practice}
A circle has area $50\pi\;\text{cm}^2$. Find its circumference.
\end{practice}
\begin{exercise}
A cylindrical grain silo has radius 2~m and height 6~m. Find its volume and
total surface area. Interpret each quantity in the context of grain storage.
\end{exercise}
\end{exercises}

% ============================================================
% PART III --- TRIGONOMETRY
% ============================================================
\part{Trigonometry: Cycles and Rotational Systems}

\chapter{Angles and Circular Geometry}

\section{The Circle as Generator of Cycles}

A circle of radius $r$ centred at the origin is the solution set of
$x^2+y^2=r^2$. As a point moves around the circle, its $x$- and
$y$-coordinates each trace out a periodic oscillation. Trigonometry is the
mathematics of this periodicity.

\begin{history}[Cyclic patterns in astrophysics]
Cecilia Payne-Gaposchkin also studied Cepheid variable stars, whose
brightness oscillates with remarkable regularity. The period of oscillation
encodes the star's intrinsic luminosity --- a fact used to measure cosmic
distances. Trigonometric modelling of such cycles was central to her work
\citep{PayneGaposchkin1925}.
\end{history}

\begin{exercises}
\begin{practice}
State the arc length formula $s=r\theta$ (where $\theta$ is in radians).
Find the arc length subtended by an angle of $\pi/3$ on a circle of radius 5~cm.
\end{practice}
\end{exercises}

\chapter{Sine, Cosine, and the Unit Circle}

\section{The Fundamental Definitions}

\begin{mydefinition}[Sine and Cosine]
For a point on the unit circle at angle $\theta$ from the positive $x$-axis,
\[
  \cos\theta = x\text{-coordinate},\qquad
  \sin\theta = y\text{-coordinate}.
\]
Equivalently, in a right triangle with adjacent side $a$, opposite
side $o$, and hypotenuse $h$:
\[
  \cos\theta=\frac{a}{h},\qquad\sin\theta=\frac{o}{h},\qquad\tan\theta=\frac{o}{a}.
\]
\end{mydefinition}

\begin{theorem}[Pythagorean Identity]
For all $\theta$,
\[
  \sin^2\theta+\cos^2\theta=1.
\]
The squared components of a unit vector always sum to one.
\end{theorem}

\begin{proofsketch}
A point $(\cos\theta,\sin\theta)$ lies on the unit circle by definition.
The circle's equation $x^2+y^2=1$ gives the identity directly.
\end{proofsketch}

\section{Exact Values}

\begin{center}
\renewcommand{\arraystretch}{1.55}
\begin{tabular}{ccccc}
\toprule
$\theta$ & radians & $\sin\theta$ & $\cos\theta$ & $\tan\theta$ \\
\midrule
$0°$   & $0$        & $0$            & $1$            & $0$ \\
$30°$  & $\pi/6$    & $1/2$          & $\sqrt{3}/2$   & $1/\sqrt{3}$ \\
$45°$  & $\pi/4$    & $\sqrt{2}/2$   & $\sqrt{2}/2$   & $1$ \\
$60°$  & $\pi/3$    & $\sqrt{3}/2$   & $1/2$          & $\sqrt{3}$ \\
$90°$  & $\pi/2$    & $1$            & $0$            & undef \\
\bottomrule
\end{tabular}
\end{center}

\begin{exercises}
\begin{practice}
Without a calculator, compute $\sin(120°)$ and $\cos(135°)$.
\end{practice}
\begin{exercise}
Prove that $\sin(90°-\theta)=\cos\theta$ using the unit circle definitions.
\end{exercise}
\begin{challenge}
The rotation matrix $R(\theta)$ preserves the norm of every vector. Express
this statement as an identity involving $\sin\theta$ and $\cos\theta$, and
relate it to the Pythagorean identity.
\end{challenge}
\end{exercises}

\chapter{Trigonometric Functions and Graphs}

\section{Periodic Functions}

\begin{mydefinition}[Period]
A function $f$ is \textbf{periodic} with period $T>0$ if $f(t+T)=f(t)$ for
all $t$. The sine and cosine functions have period $2\pi$.
\end{mydefinition}

The general sinusoidal function
\[
  f(t)=A\sin(\omega t+\phi)+D
\]
has \textbf{amplitude} $A$, \textbf{angular frequency} $\omega$,
\textbf{phase} $\phi$, and \textbf{vertical shift} $D$. The period is
$T=2\pi/\omega$.

\begin{center}
\begin{tikzpicture}
\begin{axis}[
  width=10cm, height=4cm,
  xlabel={$t$}, ylabel={$f(t)$},
  xmin=0, xmax=4*pi, ymin=-1.6, ymax=1.6,
  xtick={0,pi,2*pi,3*pi,4*pi},
  xticklabels={$0$,$\pi$,$2\pi$,$3\pi$,$4\pi$},
  axis lines=center, grid=both, grid style={ctGray!20},
  tick label style={font=\small},
]
  \addplot[ctBlue,thick,smooth,samples=120,domain=0:4*pi]{sin(deg(x))};
  \addplot[ctTeal,thick,smooth,dashed,samples=120,domain=0:4*pi]{cos(deg(x))};
  \legend{$\sin t$, $\cos t$}
\end{axis}
\end{tikzpicture}
\end{center}

\section{Variable Stars: Trigonometric Modelling}

\begin{myexample}[Payne-Gaposchkin: stellar brightness cycles]
A Cepheid variable star's brightness is modelled by
\[
  B(t)=10+2\sin\!\left(\frac{2\pi}{5}t\right),
\]
with amplitude 2 and equilibrium brightness $=10$.

At $t=2$ days:
\[
  B(2)=10+2\sin\!\left(\frac{4\pi}{5}\right)
      =10+2\sin(144°)\approx10+2(0.588)=11.18.
\]
This model was used by Payne-Gaposchkin to link oscillation period to
distance measurement.
\end{myexample}

\begin{exercises}
\begin{practice}
Identify the amplitude, period, and phase of $f(t)=3\sin(2t-\pi/4)$.
\end{practice}
\begin{exercise}
Gopnik \citep{Gopnik2009} notes that children alternate between exploration
and exploitation in a roughly periodic pattern. Model the exploration
intensity as $E(t)=0.5+0.5\sin(\omega t)$. When is exploration at maximum?
When is it at minimum? What is the period if the cycle is one week?
\end{exercise}
\end{exercises}

\chapter{Modeling Cyclic Phenomena}

\section{Harmonic Motion}

Simple harmonic motion arises whenever a restoring force is proportional to
displacement. The governing equation is
\[
  \frac{d^2x}{dt^2}+\omega^2 x=0,
\]
with general solution $x(t)=A\cos(\omega t)+B\sin(\omega t)$.

\begin{myexample}[Wave interference]
Two waves $f_1=A\sin(\omega t)$ and $f_2=A\sin(\omega t+\delta)$ combine to
\[
  f_1+f_2=2A\cos\!\left(\frac{\delta}{2}\right)\sin\!\left(\omega t+\frac{\delta}{2}\right).
\]
When $\delta=\pi$, amplitude vanishes (destructive interference). The constraint
$\delta=\pi$ is exactly the condition for cancellation.
\end{myexample}

\begin{exercises}
\begin{exercise}
Two sides of a triangle are 7 and 9; the included angle is $40°$. Use the
Law of Cosines to find the third side. Show all unit checks.
\end{exercise}
\begin{challenge}
Juarrero \citep{Juarrero1999} argues that organisms maintain coherent
behaviour through nested cycles of constraint. Model two interacting
oscillators $x(t)=\sin(t)$, $y(t)=\sin(t+\phi)$. For what value of $\phi$
do they cancel? For what value do they reinforce? Discuss the biological
analogy.
\end{challenge}
\end{exercises}

\chapter{Rotations and Harmonic Motion}

\section{Circular Motion and Projections}

A point moving at constant angular velocity $\omega$ on a circle of radius
$R$ has position
\[
  (x,y)=(R\cos(\omega t), R\sin(\omega t)).
\]
The $x$- and $y$-projections are independent sinusoidal oscillations, which is
why circular motion and wave motion share the same mathematics.

\begin{center}
\begin{tikzpicture}[scale=1.2]
  \draw[ctGray,thin] (-1.5,0)--(1.5,0);
  \draw[ctGray,thin] (0,-1.5)--(0,1.5);
  \draw[ctBlue!70,thick] (0,0) circle (1);
  \def\t{50}
  \coordinate (P) at ({cos(\t)},{sin(\t)});
  \draw[ctTeal,very thick] (0,0)--(P);
  \draw[ctRed,thick,dashed] ({cos(\t)},0)--(P);
  \draw[ctGold,thick,dashed] (0,0)--({cos(\t)},0);
  \filldraw[ctBlue] (P) circle (1.5pt);
  \node[font=\scriptsize,right] at (P){$P=(\cos\omega t,\sin\omega t)$};
\end{tikzpicture}
\end{center}

\begin{exercises}
\begin{exercise}
Show that if a point traces the unit circle at constant speed, its $x$- and
$y$-velocities are $-\sin(\omega t)$ and $\cos(\omega t)$ respectively.
Verify using the Pythagorean identity that the speed is constant.
\end{exercise}
\end{exercises}

% ============================================================
% PART IV --- STOICHIOMETRY
% ============================================================
\part{Stoichiometry: Conservation and Transformation in Chemistry}

\chapter{Atoms, Molecules, and Chemical Formulas}

\section{The Atomic Hypothesis and Conservation}

In ordinary chemical reactions, atoms are neither created nor destroyed ---
they are rearranged. This is the \textbf{law of conservation of mass}, the
chemical analogue of every balance principle studied so far.

\begin{mydefinition}[Mole and Molar Mass]
The \textbf{mole} (mol) is the SI unit of chemical amount: one mole contains
$N_A=6.022\times10^{23}$ particles (Avogadro's number). The
\textbf{molar mass} $M$ (g/mol) is the mass of one mole of a substance.
The relationship
\[
  n = \frac{m}{M}
\]
converts between mass and amount of substance.
\end{mydefinition}

\begin{history}[Payne-Gaposchkin: hydrogen and helium in stars]
Payne-Gaposchkin's 1925 discovery that stars are overwhelmingly hydrogen and
helium was resisted because it contradicted the then-prevailing view. Her
argument was essentially stoichiometric: matching observed spectral absorption
strengths to theoretical abundances. The molar proportions of H to He in a
star constrain its entire energy budget \citep{PayneGaposchkin1925}.
\end{history}

\begin{myexample}[Moles of water]
How many moles are present in 36~g of $\mathrm{H_2O}$?
\[
  M(\mathrm{H_2O})=2(1.008)+16.00=18.016\gpm.
  \qquad
  n=\frac{36}{18.016}\approx2.0\mol.
\]
\end{myexample}

\begin{exercises}
\begin{practice}
Find the molar mass of (a)~NaCl; (b)~CO$_2$; (c)~C$_6$H$_{12}$O$_6$.
\end{practice}
\begin{exercise}
How many molecules are in 44~g of CO$_2$? Express in scientific notation.
\end{exercise}
\begin{challenge}
Payne-Gaposchkin computed that the Sun's atmosphere is approximately 92\%
hydrogen by number of atoms. If the Sun's atmosphere contains $N$ total atoms,
write an expression for the number of hydrogen atoms and the number of
helium atoms (assuming the remaining 8\% is helium). What mole ratio does
this represent?
\end{challenge}
\end{exercises}

\chapter{Balancing Chemical Equations}

\section{Conservation as a Linear System}

\begin{theorem}[Stoichiometric Conservation]
A chemical equation is balanced if and only if, for each element, the total
atom count on the left side equals the total atom count on the right side.
Equivalently, the stoichiometric coefficient vector $\gamma$ satisfies
$N\gamma=\mathbf{0}$, where $N$ is the element--species incidence matrix.
\end{theorem}

\begin{myexample}[Combustion of ethane]
Balance $\mathrm{C_2H_6+O_2\to CO_2+H_2O}$.

Assign coefficients $a,b,c,d$ and write conservation equations:
\begin{align*}
  \text{C}: & \quad 2a=c \\
  \text{H}: & \quad 6a=2d \\
  \text{O}: & \quad 2b=2c+d
\end{align*}
Set $a=1$: $c=2$, $d=3$, $b=7/2$. Multiply by 2:
\[
  2\mathrm{C_2H_6}+7\mathrm{O_2}\to4\mathrm{CO_2}+6\mathrm{H_2O}.
\]
\end{myexample}

\begin{center}
\begin{tikzpicture}[node distance=1.4cm, scale=0.9]
  % Reaction network
  \node[qty,fill=ctLightGold] (CH4) {CH$_4$};
  \node[qty,fill=ctLightBlue,right=2.2cm of CH4] (CO2) {CO$_2$};
  \node[qty,fill=ctLightTeal,below=of CH4] (O2) {O$_2$};
  \node[qty,fill=ctLightTeal,below=of CO2] (H2O) {H$_2$O};
  \draw[arrow] (CH4) -- (CO2) node[midway,above,font=\scriptsize]{reaction};
  \draw[arrow] (O2) -- (CO2);
  \draw[arrow] (O2) -- (H2O);
  \draw[arrow] (CH4) -- (H2O);
  \node[font=\scriptsize,ctGray,below=1.8cm of O2]{$\mathrm{CH_4+2O_2\to CO_2+2H_2O}$};
\end{tikzpicture}
\end{center}

\begin{exercises}
\begin{practice}
Balance (a)~$\mathrm{Fe+O_2\to Fe_2O_3}$;\;
(b)~$\mathrm{C_3H_8+O_2\to CO_2+H_2O}$.
\end{practice}
\begin{exercise}
Write the stoichiometric matrix $N$ for the reaction
$\mathrm{CH_4+2O_2\to CO_2+2H_2O}$ and verify $N\gamma=\mathbf{0}$.
\end{exercise}
\begin{challenge}
Balance $\mathrm{KMnO_4+HCl\to KCl+MnCl_2+H_2O+Cl_2}$ by writing and
solving the linear system. This is a non-trivial redox reaction with 7
unknowns.
\end{challenge}
\end{exercises}

\chapter{Moles and Quantitative Chemistry}

\section{Stoichiometric Calculations}

\begin{myexample}[Mass of water produced]
From $2\mathrm{H_2+O_2\to2H_2O}$, how many grams of water form from 4~g of
$\mathrm{H_2}$ (excess $\mathrm{O_2}$)?
\begin{align*}
  n(\mathrm{H_2}) &= \frac{4\;\text{g}}{2.016\gpm} \approx1.984\mol\\
  n(\mathrm{H_2O}) &= 1.984\mol\times\frac{2\;\text{mol H}_2\text{O}}{2\;\text{mol H}_2}
    = 1.984\mol\\
  m(\mathrm{H_2O}) &= 1.984\times18.016 \approx 35.7\;\text{g}
\end{align*}
\end{myexample}

\begin{exercises}
\begin{practice}
From $\mathrm{C_3H_8+5O_2\to3CO_2+4H_2O}$, how many grams of CO$_2$ are
produced by burning 11~g of propane ($M=44.1\gpm$)?
\end{practice}
\end{exercises}

\chapter{Reaction Ratios and Limiting Reagents}

\section{The Limiting Reagent}

\begin{mydefinition}[Limiting Reagent]
The \textbf{limiting reagent} is the reactant consumed first, setting the
maximum possible yield. All other reactants are \textbf{excess reagents}.
\end{mydefinition}

\begin{myexample}[Juarrero: constraint propagation in a reaction]
Consider $\mathrm{N_2+3H_2\to2NH_3}$. Begin with 2~mol $\mathrm{N_2}$ and
5~mol $\mathrm{H_2}$.

$\mathrm{N_2}$ alone could produce $4\;\text{mol NH}_3$; $\mathrm{H_2}$ alone
could produce $5\times\tfrac{2}{3}\approx3.33\;\text{mol NH}_3$.

$\mathrm{H_2}$ is limiting. Yield $=3.33\;\text{mol NH}_3$.

Juarrero \citep{Juarrero1999} calls this \textit{constraint cascade}: the
scarcest resource propagates its limitation through the entire network.
\end{myexample}

\section{Percent Yield}

\[
  \%\text{yield}
  = \frac{\text{actual yield}}{\text{theoretical yield}}\times100\%.
\]

\begin{exercises}
\begin{exercise}
How many moles of excess $\mathrm{N_2}$ remain after the reaction in the
example above?
\end{exercise}
\begin{challenge}
Model a three-step reaction cascade $A+B\to C$, $C+D\to E$, $E+F\to G$.
Starting with 4~mol each of $A,B,D,F$ and none of the intermediates,
determine the theoretical yield of $G$. Identify which reagent is limiting
at each stage.
\end{challenge}
\end{exercises}

\chapter{Reaction Networks and Conservation Laws}

\section{Dynamic Stoichiometry}

In a reaction network, concentrations $\mathbf{c}$ evolve by
\[
  \frac{d\mathbf{c}}{dt} = S\,\mathbf{r}(\mathbf{c}),
\]
where $S$ is the stoichiometric matrix and $\mathbf{r}(\mathbf{c})$ is the
vector of reaction rates. Conservation is encoded in $\ker(S)$.

\begin{myexample}[Continuity in disguise]
The equation above is the discrete analogue of the field-theoretic continuity
equation
\[
  \pd{\rho}{t} + \nabla\cdot(\rho\mathbf{v}) = 0,
\]
describing conservation of density $\rho$ in a flowing medium with velocity
$\mathbf{v}$. Algebra, stoichiometry, and field theory all encode the same
structural idea: \emph{what flows in must flow out}.
\end{myexample}

\begin{deepnote}
Appendix~E derives the RSVP field equations as the continuum limit of this
conservation structure.
\end{deepnote}

% ============================================================
% PART V --- INTEGRATION
% ============================================================
\part{Integration: Mathematics of Structured Systems}

\chapter{Constraint Networks}

\section{Networks as Mathematical Objects}

A \textbf{network} (or graph) consists of \textbf{nodes} (vertices) and
\textbf{edges} (links). When nodes represent quantities and edges represent
constraints, the network encodes the structure of a system.

\begin{myexample}[Social constraint network after Arendt]
Arendt's analysis of political space can be formalised as a constraint network
in which nodes are actors and directed edges are communicative acts
\citep{Arendt1958}. The system is stable when the network satisfies certain
consistency conditions --- analogous to the solvability of a linear system.
\end{myexample}

\begin{myexample}[Ecological resource balance after Le Guin]
The allocation system $F=2L$, $E=F+L$ from Chapter~4 is a two-node constraint
network. Le Guin's utopian societies must maintain exactly such networks at
equilibrium \citep{LeGuin1974}.
\end{myexample}

\chapter{Cycles, Flows, and Equilibrium}

\section{Three Kinds of Balance}

All the subjects in this book converge on three structural ideas:

\begin{center}
\renewcommand{\arraystretch}{1.6}
\begin{tabular}{>{\bfseries}p{3cm} p{4cm} p{5cm}}
\toprule
\color{ctBlue}Type & \color{ctBlue}Domain & \color{ctBlue}Equation \\
\midrule
Algebraic balance  & variable equations     & $ax=b$ \\
Spatial balance    & Pythagorean geometry   & $a^2+b^2=c^2$ \\
Rotational balance & unit circle            & $\sin^2+\cos^2=1$ \\
Material balance   & stoichiometry          & $N\gamma=\mathbf{0}$ \\
Dynamic balance    & field continuity       & $\partial_t\rho+\nabla\cdot(\rho\mathbf{v})=0$ \\
\bottomrule
\end{tabular}
\end{center}

\section{Equilibrium as Fixed Point}

\begin{mydefinition}[Equilibrium]
A system state $\mathbf{x}^*$ is an \textbf{equilibrium} if applying the
system's dynamics to $\mathbf{x}^*$ returns $\mathbf{x}^*$:
$T(\mathbf{x}^*)=\mathbf{x}^*$.
\end{mydefinition}

\chapter{Scaling and Multilevel Structure}

\section{Patterns That Recur Across Scales}

\begin{theorem}[Self-Similarity Under Scaling]
If a relationship $y=kx^p$ holds at one scale, and the same exponent $p$
governs the relationship at larger or smaller scales, the pattern is
\textbf{scale-invariant}.
\end{theorem}

The inverse-square law $I=L/(4\pi r^2)$ is scale-invariant: it holds whether
$r$ is one metre or one light-year.

\begin{myexample}[Hazen: mineral evolution]
Hazen et al.\ \citep{Hazen2008} trace how the diversity of mineral species on
Earth grew in proportion to the complexity of planetary processes. At each
scale --- atomic bonding, crystal structure, tectonic activity, biological
metabolism --- the same constraint logic applies: new mineral species form
precisely when new constraints (temperature, pressure, chemical environment)
are satisfied.
\end{myexample}

\chapter{From Elementary Systems to Scientific Models}

\section{The Architecture of Science}

The tools introduced in this book form the foundation of all quantitative
science. Every advanced model is an elaboration of the same constraint
structure encountered in Chapter~1.

\begin{center}
\begin{tikzpicture}[node distance=1.2cm and 2.8cm, scale=0.9]
  \node[qty,fill=ctLightBlue] (A) {Algebra};
  \node[qty,fill=ctLightBlue,right=of A] (G) {Geometry};
  \node[qty,fill=ctLightTeal,below=of A] (T) {Trig};
  \node[qty,fill=ctLightTeal,right=of T] (S) {Stoich};
  \node[qty,fill=ctLightGold,below=2cm of T] (F) {Fields};
  \node[qty,fill=ctLightGold,right=of F] (C) {Cats};
  \draw[eqedge] (A)--(G); \draw[eqedge] (A)--(T); \draw[eqedge] (A)--(S);
  \draw[eqedge] (G)--(T); \draw[eqedge] (G)--(S);
  \draw[eqedge] (T)--(F); \draw[eqedge] (S)--(F);
  \draw[eqedge] (F)--(C); \draw[eqedge] (A)--(F);
\end{tikzpicture}
\end{center}

\begin{principle}
Scientific understanding grows by identifying patterns of constraint and
transformation across domains. Transformation reveals structure; invariance
reveals law.
\end{principle}

% ============================================================
% APPENDICES
% ============================================================
\appendix

\chapter{Algebra as Constraint Before Content}

\section{The General Constraint System}

Let $x_1,\dots,x_n\in\R$ be system variables. A \textbf{constraint system}
is a collection of equations
\[
  F_\alpha(x_1,\dots,x_n)=0,\quad\alpha=1,\dots,m.
\]
The solution locus is $\mathcal{C}=\{x\in\R^n\mid F_\alpha(x)=0\;\forall\alpha\}$.

For linear constraints, $F_\alpha(x)=A_\alpha\cdot x - b_\alpha$, so the
system becomes $A\mathbf{x}=\mathbf{b}$ with $A\in\R^{m\times n}$.

\section{Admissibility and Relaxation}

When $m>n$ the system is overdetermined and may be inconsistent. Define the
\textbf{residual functional}
\[
  \mathcal{R}(\mathbf{x})=\sum_{\alpha=1}^m|F_\alpha(\mathbf{x})|^2.
\]
Approximate solutions minimise $\mathcal{R}$. This is the \textbf{least-squares
principle}, which is also the precursor to entropic relaxation in the RSVP
framework (Appendix~E).

\begin{proposition}[Least Squares]
The minimiser of $\mathcal{R}(\mathbf{x})=\norm{A\mathbf{x}-\mathbf{b}}^2$
satisfies the \textbf{normal equations} $A^\top A\mathbf{x}=A^\top\mathbf{b}$.
\end{proposition}

\section{Degrees of Freedom}

When $m<n$, the system is underdetermined and the solution locus is a
subspace of dimension $n-\mathrm{rank}(A)$. These degrees of freedom
represent the system's capacity to vary while remaining consistent ---
a structural feature central to RSVP field configurations.

\chapter{Geometry as Structured Invariance}

\section{Euclidean Space and the Inner Product}

On $\R^n$, the standard inner product is
$\langle x,y\rangle=\sum_{i=1}^n x_iy_i$,
with induced norm $\norm{x}=\sqrt{\langle x,x\rangle}$.

The distance formula and Pythagorean theorem both follow:
$\norm{x-y}^2=\norm{x}^2-2\langle x,y\rangle+\norm{y}^2$.

\section{Configuration Spaces}

The RSVP theory works not over finite-dimensional $\R^n$ but over an
infinite-dimensional configuration space:
\[
  \mathcal{M}=\left\{(\Phi,\mathbf{v},S)\mid
  \Phi:\Omega\to\R,\;\mathbf{v}:\Omega\to\R^3,\;S:\Omega\to\R\right\}.
\]
An $L^2$-type inner product may be defined as
\[
  \langle X_1,X_2\rangle=
  \int_\Omega\!\left(\Phi_1\Phi_2+\mathbf{v}_1\cdot\mathbf{v}_2+S_1S_2\right)dx.
\]
The Euclidean geometry of finite dimensions extends naturally to this
infinite-dimensional setting.

\chapter{Trigonometry as Oscillation, Rotation, and Recurrence}

\section{The Rotation Group}

Define $R(\theta)\in\R^{2\times2}$ as in Chapter~9. These matrices form the
group $\mathrm{SO}(2)$ under composition, with $R(\theta_1)R(\theta_2)=R(\theta_1+\theta_2)$.

The Pythagorean identity $\sin^2\theta+\cos^2\theta=1$ is precisely the
condition $\det R(\theta)=1$.

\section{Harmonic Analysis}

The differential equation $\ddot{x}+\omega^2x=0$ governs simple harmonic
motion. Its complex solutions are $e^{\pm i\omega t}$, where
$e^{i\omega t}=\cos(\omega t)+i\sin(\omega t)$ (Euler's formula).

Linearising the RSVP system around a uniform background state yields modes
of the form $e^{i(\mathbf{k}\cdot\mathbf{x}-\omega t)}$ --- the same
trigonometric recurrence, now in three spatial dimensions.

\chapter{Stoichiometry as Conservation Algebra}

\section{The Stoichiometric Null Space}

For a reaction network with stoichiometric matrix $S\in\Z^{e\times r}$
(elements $\times$ reactions), conservation is the condition
$S\gamma=\mathbf{0}$. The admissible coefficient vectors form the
\textbf{null space} $\ker(S)$.

\section{Dynamic Stoichiometry}

Concentrations $\mathbf{c}(t)\in\R^r$ evolve by
\[
  \dot{\mathbf{c}}=S\,\mathbf{r}(\mathbf{c}),
\]
the \emph{Chemical Master Equation} in its deterministic (mass-action) limit.

The connection to field theory: replacing $\mathbf{c}$ by a density field
$\rho(x,t)$ and the finite sum by a spatial integral yields the continuity
equation $\partial_t\rho+\nabla\cdot(\rho\mathbf{v})=0$.

\chapter{RSVP Field Theory: Foundational Derivation}

\section{The Fundamental Fields}

Let $\Omega\subset\R^3$ be a spatial domain. The RSVP theory posits three
coupled fields:
\[
  \Phi(x,t)\in\R,\quad
  \mathbf{v}(x,t)\in\R^3,\quad
  S(x,t)\in\R,
\]
representing plenum density, vector flow, and entropy density respectively.

\section{Governing Equations}

\begin{align}
  \partial_t\Phi+\nabla\cdot(\Phi\mathbf{v})
  &= D_\Phi\Delta\Phi+F_\Phi(\Phi,\mathbf{v},S), \tag{A}\\
  \partial_t\mathbf{v}+(\mathbf{v}\cdot\nabla)\mathbf{v}
  &= -\nabla P(\Phi,S)+D_v\Delta\mathbf{v}+F_v(\Phi,\mathbf{v},S), \tag{B}\\
  \partial_t S+\nabla\cdot(S\mathbf{v})
  &= D_S\Delta S+\sigma(\Phi,\mathbf{v},S). \tag{C}
\end{align}

Each left-hand side is an advection--transport operator; the diffusion terms
smooth spatial gradients; the source terms encode nonlinear coupling.

\section{Lyapunov Functional and Entropic Descent}

Define
\[
  \mathcal{F}[\Phi,\mathbf{v},S]
  =\int_\Omega\!\left[
    \tfrac12|\mathbf{v}|^2
    +U(\Phi,S)
    +\tfrac{\kappa_\Phi}{2}|\nabla\Phi|^2
    +\tfrac{\kappa_S}{2}|\nabla S|^2
  \right]dx.
\]
Dissipative dynamics descend on $\mathcal{F}$:
$\partial_t X = -\frac{\delta\mathcal{F}}{\delta X}$.
This makes precise the intuition that ``space falls outward'' toward lower
free energy --- the entropic smoothing picture underlying RSVP cosmology.

\section{Linear Perturbation Theory}

Expand around a homogeneous state
$\Phi=\Phi_0+\delta\Phi$, $\mathbf{v}=\delta\mathbf{v}$, $S=S_0+\delta S$.
The linearised system admits Fourier modes
$\delta X\sim e^{i(\mathbf{k}\cdot x-\omega t)}$,
whose dispersion relation $\omega(\mathbf{k})$ determines stability and
propagation speed.

\chapter{Entropic Redshift and Non-Expansion Cosmology}

\section{Redshift Without Metric Expansion}

Let $\mathcal{E}(x,t)$ be an entropic attenuation field. Photon frequency
along ray $x(\lambda)$ evolves as
\[
  \frac{d\nu}{d\lambda}=-\alpha\,\mathcal{E}(x(\lambda),t)\,\nu,
\]
with solution
\[
  \nu(\lambda)=\nu_0\exp\!\left(-\alpha\int_0^\lambda\mathcal{E}\,d\lambda'\right).
\]
The redshift is
$1+z = \exp\!\left(\alpha\int_0^\lambda\mathcal{E}\,d\lambda'\right)$.
For weak, homogeneous $\mathcal{E}$:
\[
  z\approx\alpha\mathcal{E}\cdot\lambda,
\]
recovering the observed linear distance--redshift law --- without assuming
metric expansion of space.

\chapter{TARTAN: Recursive Tiling and Annotated Noise}

\section{Multiscale Partition}

Let $\Omega$ admit a nested sequence of partitions
$\mathcal{T}_0\supseteq\mathcal{T}_1\supseteq\cdots$ with refinement maps
$\pi_k:\mathcal{T}_{k+1}\to\mathcal{T}_k$.

\section{Annotated Field Updates}

Each tile $T\in\mathcal{T}_k$ carries state
$X_k(T)=(\Phi_k,\mathbf{v}_k,S_k,\eta_k)$,
where $\eta_k$ is a semantic / trajectory annotation.

Recursive update:
\[
  X_{k+1}(T')=F_k\!\left(X_k(\pi_k(T')),\,\xi_k(T')\right),
\]
where noise decomposes as $\xi=\xi_\text{white}+\xi_\text{structured}$,
the structured part carrying provenance metadata.

History dependence is encoded by a trajectory buffer:
\[
  \mathcal{B}(T,t)=\{X(T,t-\tau_j)\}_{j=1}^L,\qquad
  X(T,t+\Delta t)=G\!\left(X(T,t),\mathcal{B}(T,t),\xi(T,t)\right).
\]

\chapter{Simulated Agency as Sparse Projection}

\section{The Projection Operator}

Let $\mathcal{W}$ be the world state space and $\mathcal{M}$ the internal
model space. An agent is characterised by a projection
$P:\mathcal{W}\to\mathcal{M}$, $m=P(w)$.

Optimal projection minimises
\[
  m^* = \arg\min_m\bigl[L_\text{task}(m;w)+\beta\,C(m)\bigr],
\]
trading off task loss against representational complexity.

\section{The Agency Loop}

\[
  w_{t+1}=F(w_t,a_t),\qquad
  m_t=P(w_t),\qquad
  a_t=\pi(m_t).
\]

This recursive loop connects world, internal model, and action. It is a
dynamical system on $\mathcal{W}\times\mathcal{M}$. Equilibria of the loop
correspond to stable behaviours; periodic orbits correspond to habitual
cycles. The geometric analysis of Chapter~9 applies directly.

\chapter{Spherepop and Event-History Formalism}

\section{Events as Primitives}

Let $E$ be a set of events with causal order $\preceq$: $e_i\preceq e_j$
means $e_i$ is causally prior to $e_j$.

A \textbf{state} is derived:
$\sigma_t=\Sigma(E_{\le t})$, where $\Sigma$ is a summarisation operator.

\begin{proposition}
Ordinary state-update equations $x_{t+1}=f(x_t)$ arise as compressed
approximations to richer event-structured evolution in which $E_{\le t}$
is the fundamental datum.
\end{proposition}

\section{Admissible Futures}

\[
  \mathrm{Poss}(E_{\le t})
  =\{e\in E\mid E_{\le t}\cup\{e\}\text{ is causally consistent}\}.
\]

This turns history into a \emph{generator} of possible futures. The
``append-only algebra'' of Spherepop is precisely the monotone extension
$E_{\le t}\mapsto E_{\le t}\cup\{e\}$ for $e\in\mathrm{Poss}(E_{\le t})$.

\chapter{Category-Theoretic Semantics of Constraint Systems}

\section{Categories}

A \textbf{category} $\mathcal{C}$ consists of:
\begin{itemize}
\item objects $\Ob(\mathcal{C})$;
\item for each pair $A,B$, a set $\Hom(A,B)$ of morphisms;
\item an associative composition law $\circ$ and identity morphisms $\id_A$.
\end{itemize}

\textbf{Examples:} $\mathbf{Set}$ (sets, functions); $\mathbf{Vect}_k$
(vector spaces, linear maps); $\mathbf{RxnNet}$ (species configurations,
stoichiometric maps).

\section{Symmetric Monoidal Structure}

A \textbf{symmetric monoidal category} $(\mathcal{C},\otimes,I)$ supports
parallel composition. Semantic modules are objects; consistency-preserving
transformations are morphisms; simultaneous operation is $\otimes$.

A merge operation, when it exists, is a \textbf{colimit}: the universal
object receiving compatible maps from all merged components.

\chapter{Sheaf-Theoretic Gluing and Local-to-Global Consistency}

\section{Presheaves and Sheaves}

Let $X$ be a topological space (the context space). A \textbf{presheaf}
$\mathcal{F}$ assigns data $\mathcal{F}(U)$ to each open set $U$, with
restriction maps $\rho_{UV}:\mathcal{F}(U)\to\mathcal{F}(V)$ for $V\subseteq U$.

$\mathcal{F}$ is a \textbf{sheaf} if: whenever sections
$\{s_i\in\mathcal{F}(U_i)\}$ satisfy $\rho_{U_i,U_i\cap U_j}(s_i)=\rho_{U_j,U_i\cap U_j}(s_j)$
on all overlaps, there exists a unique global section $s\in\mathcal{F}(\bigcup U_i)$
restricting to each $s_i$.

\begin{principle}
Sheaf coherence is the precise mathematical form of local-to-global
consistency. Every chapter of this book has been practicing local-to-global
reasoning in disguise: from triangle ratios to stellar spectra, from
stoichiometric balance to network equilibrium.
\end{principle}

\chapter{Variational Principles and Entropic Descent}

\section{The Euler-Lagrange Equation}

Let $\mathcal{J}[u]=\int_\Omega\mathcal{L}(u,\nabla u,x)\,dx$. Critical
points satisfy the \textbf{Euler-Lagrange equation}:
\[
  \pd{\mathcal{L}}{u}-\nabla\cdot\left(\pd{\mathcal{L}}{(\nabla u)}\right)=0.
\]

\section{Gradient Flow}

Dissipative dynamics minimising $\mathcal{J}$ satisfy:
\[
  \partial_t u = -\frac{\delta\mathcal{J}}{\delta u}.
\]

This single equation subsumes: algebraic residual minimisation
($\mathcal{J}=\norm{A\mathbf{x}-\mathbf{b}}^2$), geometric shortest paths,
chemical equilibration, and entropic smoothing in RSVP --- all are gradient
flows on functionals appropriate to their domain.

\chapter{Worked Derivations: Bridges Between the Introductory Subjects and the Theory}

\section{From Linear Algebra to Stoichiometry to Field Theory}

\textit{Step 1.} Balancing $2\mathrm{H_2+O_2\to2H_2O}$ is solving
$S\gamma=\mathbf{0}$ with
$S=\begin{pmatrix}2&0&2\\0&2&1\end{pmatrix}$.

\textit{Step 2.} Dynamic generalisation: $\dot{\mathbf{c}}=S\mathbf{r}(\mathbf{c})$.

\textit{Step 3.} Continuum limit: replace $c_i(t)$ by $\rho(x,t)$ and
the discrete sum by a spatial divergence, yielding
$\partial_t\rho+\nabla\cdot(\rho\mathbf{v})=0$.

Atom counting and field theory are the same idea at different scales.

\section{From Distance Formula to $L^2$ Metric}

\textit{Step 1.} Euclidean distance $d(P,Q)=\norm{P-Q}$ in $\R^n$.

\textit{Step 2.} Extend to function space:
$d(f,g)=\norm{f-g}_{L^2}=\left(\int_\Omega|f-g|^2dx\right)^{1/2}$.

\textit{Step 3.} RSVP uses exactly this $L^2$ metric on field configurations
to define entropic distance between system states.

\section{From Trigonometric Oscillation to RSVP Eigenmodes}

\textit{Step 1.} $x(t)=A\cos(\omega t)$ solves $\ddot{x}+\omega^2x=0$.

\textit{Step 2.} Linearise RSVP around uniform background. The perturbation
$\delta\Phi$ satisfies $\partial_t^2(\delta\Phi)+\omega^2(\mathbf{k})\delta\Phi=0$
for each Fourier mode $\mathbf{k}$.

\textit{Step 3.} Each mode oscillates sinusoidally; the dispersion relation
$\omega(\mathbf{k})$ encodes how the RSVP field propagates information.

% ============================================================
% BACK MATTER
% ============================================================
\backmatter

% ---- Glossary -----------------------------------------------
\chapter*{Glossary}
\addcontentsline{toc}{chapter}{Glossary}

\begin{multicols}{2}\small
\begin{description}[labelindent=0pt,leftmargin=0.8em,style=nextline]
\item[Amplitude] maximum displacement of a periodic function from its midline.
\item[Avogadro's number] $N_A=6.022\times10^{23}$; particles per mole.
\item[Constraint] a condition restricting admissible values of variables.
\item[Conservation law] a relation asserting that a quantity is preserved under transformation.
\item[Dimensional analysis] tracking units through calculations.
\item[Equilibrium] a fixed point of a dynamical system.
\item[Event-history] a causally ordered sequence of occurrences; primitive in the Spherepop formalism.
\item[Function] a rule assigning exactly one output to each input.
\item[Gradient flow] dynamics $\partial_t u=-\delta\mathcal{J}/\delta u$ descending a functional.
\item[Identity (trig.)] an equation true for all valid angles.
\item[Invariant] a quantity unchanged by a transformation.
\item[Isometry] a transformation preserving distances.
\item[Limiting reagent] the reactant consumed first in a reaction.
\item[Linear system] a system of equations of the form $A\mathbf{x}=\mathbf{b}$.
\item[Molar mass] mass per mole of a substance (g/mol).
\item[Mole] SI unit of chemical amount.
\item[Null space] $\ker(A)=\{\mathbf{x}\mid A\mathbf{x}=\mathbf{0}\}$.
\item[Period] smallest positive interval of repetition of a periodic function.
\item[Pythagorean theorem] $a^2+b^2=c^2$ for a right triangle.
\item[Radian] angle subtended by arc of length equal to radius.
\item[RSVP] Relativistic Scalar-Vector Plenum field theory.
\item[Sheaf] a presheaf satisfying the gluing axiom for local data.
\item[Similarity] geometric relation in which figures differ only by scale.
\item[Slope] rate of change of a linear function.
\item[Stoichiometric matrix] element-species incidence matrix $S$ encoding conservation.
\item[TARTAN] Recursive Tiling and Annotated Noise simulation framework.
\item[Unit circle] circle of radius 1; foundation of trigonometric definitions.
\item[Variable] a letter representing an unknown or varying quantity.
\item[Variational functional] a function on function spaces, $\mathcal{J}[u]$.
\end{description}
\end{multicols}

% ---- Selected answers ---------------------------------------
\chapter*{Selected Answers}
\addcontentsline{toc}{chapter}{Selected Answers}

\textbf{Chapter 2 (Equations and Balance).}
Practice 2.1: (a) $x=5$; (b) $x=15$; (c) $x=5$; (d) $x=7$.

\textbf{Chapter 7 (Pythagorean Theorem).}
Practice 7.1: diagonal $=\sqrt{49+576}=25$~cm.

\textbf{Chapter 10 (Atoms and Moles).}
Practice 10.1: (a) 58.44; (b) 44.01; (c) 180.16~g/mol.

\textbf{Chapter 17 (Balancing).}
Practice 17.1: (a) $4\mathrm{Fe}+3\mathrm{O_2}\to2\mathrm{Fe_2O_3}$;
(b) $\mathrm{C_3H_8}+5\mathrm{O_2}\to3\mathrm{CO_2}+4\mathrm{H_2O}$.

\textbf{Chapter 18 (Moles).}
Practice 18.1: $\approx33.6$~g CO$_2$.

% ---- Abbreviated Periodic Table ------------------------------
\chapter*{Abbreviated Periodic Table}
\addcontentsline{toc}{chapter}{Abbreviated Periodic Table}

\begin{center}
\renewcommand{\arraystretch}{1.45}
\begin{tabular}{clcc}
\toprule
\textbf{Symbol} & \textbf{Name} & \textbf{$Z$} & \textbf{$M$ (g/mol)} \\
\midrule
H  & Hydrogen  & 1  & 1.008  \\
C  & Carbon    & 6  & 12.011 \\
N  & Nitrogen  & 7  & 14.007 \\
O  & Oxygen    & 8  & 15.999 \\
Na & Sodium    & 11 & 22.990 \\
Mg & Magnesium & 12 & 24.305 \\
S  & Sulfur    & 16 & 32.060 \\
Cl & Chlorine  & 17 & 35.453 \\
K  & Potassium & 19 & 39.098 \\
Ca & Calcium   & 20 & 40.078 \\
Fe & Iron      & 26 & 55.845 \\
Cu & Copper    & 29 & 63.546 \\
\bottomrule
\end{tabular}
\end{center}

% ---- Bibliography -------------------------------------------
\bibliographystyle{plainnat}
\bibliography{constraint_refs}

\end{document}
